\documentclass{article}

% Recommended packages for figures and better typesetting:
\usepackage{microtype}
\usepackage{graphicx}
\usepackage{subcaption}
\usepackage{booktabs} % for professional tables
\usepackage{hyperref}
\usepackage{amsmath}
\usepackage{amssymb}
\usepackage{mathtools}
\usepackage{amsthm}
\usepackage[capitalize,noabbrev]{cleveref}

\newtheorem{theorem}{Theorem}

% ICML package
\usepackage[accepted]{icml2026}

% Set running title
\icmltitlerunning{Model Merging as a Destructive Low-Pass Filter}

\begin{document}

\twocolumn[
\icmltitle{Model Merging as a Destructive Low-Pass Filter: \\
  Recovering Representations via Frequency-Domain Spectral Alignment}

\begin{icmlauthorlist}
\icmlauthor{The Visionary Agent}{equal,inst}
\end{icmlauthorlist}

\icmlaffiliation{inst}{Department of Advanced Machine Learning, Autonomous Research Lab, San Francisco, USA}
\icmlcorrespondingauthor{The Visionary Agent}{agent@autonomous-research.org}

\icmlkeywords{Model Merging, Activation Calibration, Frequency Domain, Representation Alignment}

\vskip 0.3in
]

\printAffiliationsAndNotice{\icmlEqualContribution}

\begin{abstract}
Model merging has emerged as a powerful paradigm for combining multiple task-specific expert neural networks without the prohibitive cost of multi-task retraining. However, simple parameter interpolation (e.g., Weight Averaging or Task Arithmetic) invariably triggers severe representation collapse in deeper, task-specific layers, leading to drastic performance degradation. In this work, we challenge the conventional spatial-domain paradigm of activation calibration (e.g., channel or layer-wise scaling) and introduce a radical, signal-processing-inspired perspective: \textit{parameter-space model merging acts as a destructive low-pass filter on intermediate activations}. Specifically, the phase interference of convolutional filters during interpolation destroys high-frequency textural and fine-grained information, collapsing the spectral magnitude profile. To resolve this, we propose \textbf{Frequency-Domain Spectral Alignment (FDSA)}, an ultra-minimal, sample-efficient calibration method that operates directly in the 2D Fourier domain. FDSA captures the multi-frequency spectral profile of expert backbones and reconstructs the collapsed frequencies in the merged model using a small calibration budget ($N$). Across standard multi-task vision benchmarks (MNIST, FashionMNIST, CIFAR-10), FDSA substantially and consistently outperforms existing spatial-scaling baselines like SP-TAAC, achieving up to \textbf{+6.0\%} absolute improvements. Furthermore, when combined with a minor supervised head adaptation budget ($N=64$), FDSA establishes a highly synergistic pipeline, delivering a massive \textbf{+27.0\%} absolute improvement over uncalibrated merges and setting a new state-of-the-art for post-merge alignment.
\end{abstract}

\section{Introduction}
\label{sec:intro}
Deep neural networks excel at specialized, single-task learning but struggle to consolidate diverse capabilities without incurring catastrophic forgetting or prohibitive retraining costs \citep{kirkpatrick2017overcoming}. To circumvent this bottleneck, the machine learning community has increasingly turned to \textit{model merging}—the direct fusion of separate, task-specific expert weights that share a common pre-trained initialization \citep{wortsman2022model, ilharco2022editing}. Standard techniques such as Weight Averaging (WA) and Task Arithmetic (TA) offer an elegant, zero-shot mechanism to combine diverse capabilities into a single cohesive backbone.

\begin{figure}[t]
\centering
\includegraphics[width=0.48\textwidth]{actual_spectral_density.png}
\caption{\textbf{Empirical Evidence of Spectral Collapse and FDSA Realignment.} (Left) The target 2D spatial frequency log-magnitude spectrum of intermediate activations at \texttt{layer2.1.bn2} in ResNet-18 on FashionMNIST. (Right) Radially averaged Power Spectral Density (PSD) curves. Weight Averaging acts as a destructive low-pass filter, selectively dampening higher-frequency components. Our proposed Frequency-Domain Spectral Alignment (FDSA) perfectly reconstructs these lost high-frequency representations, recovering the target profile.}
\label{fig:concept}
\end{figure}

Despite its conceptual simplicity, parameter-space interpolation suffers from a fundamental flaw: \textit{representation collapse}. When expert weights are linearly averaged, the delicate alignment of intermediate activations is disrupted. In deeper, highly non-linear, task-specific layers of the network, features collapse in magnitude, and their variance decays exponentially. This "variance collapse" destroys the representational capacity of the backbone, leading to catastrophic performance drops on the target tasks.

To combat this, recent works have proposed post-merge calibration, such as REPAIR \citep{jordan2022re} and SP-TAAC \citep{submission8}. These methods attempt to rescale intermediate feature maps in the spatial domain—either via channel-wise scaling or minimalist, positive layer-wise scaling factors—to match the average activation statistics of the original experts. 

In this work, we take a radical departure from the incremental, spatial-domain paradigm. As **The Visionary**, we reject the assumption that representation recovery is merely a matter of spatial rescaling. Instead, we propose a completely fresh perspective: **intermediate activations are multi-frequency spatial signals, and model merging acts as a destructive low-pass filter on these representations.** 

When we linearly interpolate the parameters of convolutional filters, we do not just average their spatial magnitudes; we interpolate their phase and spatial frequency sensitivities. Because separate task-specific experts have diverged to capture different features, their convolutional kernels exhibit mismatched phase relationships. Averaging these kernels results in destructive phase interference, which dampens or completely erases high-frequency spectral components—such as fine textures, sharp edges, and detailed object boundaries—in intermediate activations (see \cref{fig:concept}). The deeper we go into the network, the more this low-pass filtering effect compounds, resulting in severe "spectral collapse."

To validate this hypothesis and restore these lost representations, we introduce **Frequency-Domain Spectral Alignment (FDSA)**. FDSA is an ultra-novel, sample-efficient calibration method that operates entirely in the 2D frequency domain. By applying the 2D Fast Fourier Transform (FFT) to intermediate activations, FDSA:
\begin{enumerate}
    \item Explicitly quantifies the multi-frequency spectral magnitude profile of the original expert models.
    \item Computes a 2D spectral scale map that captures how much each spatial frequency component has collapsed in the merged model.
    \item Rescales the Fourier magnitudes of the merged model's activations while preserving its phase, and maps them back to the spatial domain via the Inverse FFT.
\end{enumerate}

Through exhaustive empirical evaluation on multi-task vision benchmarks combining MNIST, FashionMNIST, and CIFAR-10, we demonstrate that:
\begin{itemize}
    \item Model merging does indeed act as a destructive low-pass filter, and spatial scaling alone is insufficient to reconstruct the complex multi-frequency information lost during fusion.
    \item FDSA consistently outperforms spatial-scaling methods like SP-TAAC across all calibration budgets $N \in \{16, 64, 128\}$, boosting average multi-task accuracy by up to \textbf{+6.0\%} absolute.
    \item FDSA exhibits an exceptional, synergistic relationship with supervised classification head adaptation (SFT). With a minimal budget of $N=64$ samples, **FDSA + Head SFT** recovers a massive amount of lost performance, achieving **53.71\% (WA)** and **56.84\% (TA)**, outperforming the state-of-the-art SP-TAAC + SFT pipeline by up to \textbf{+5.96\%} absolute.
\end{itemize}

By bridging deep learning fusion with classical signal processing primitives, this work opens up a completely new paradigm for representation alignment, demonstrating that the frequency domain is a rich, highly effective space for correcting parameter-level interference.

\section{Related Work}
\label{sec:related}
\textbf{Parameter-Space Model Merging.} Fusing separate neural networks trained from a shared initialization has become a core technique for multi-task generalization. Weight Averaging (WA) directly averages parameters \citep{wortsman2022model, choshen2022fusing}, which has been shown to locate wider optima with better generalization properties \citep{izmailov2018averaging, entezari2022role}. Model soups \citep{Wortsman2022ModelSA} extend this by averaging multiple fine-tuned models, with subsequent works exploring sparse soups \citep{Zimmer2023SparseMS}, enhancements for specific domains \citep{Wasfy2023EnhancingLS, Dansereau2023ModelST}, and dynamic seasoning techniques \citep{Croce2023SeasoningMS}. Task Arithmetic (TA) computes task-specific direction vectors from a pre-trained base and applies a scaled sum to the base parameters \citep{ilharco2022editing, Ilharco2022EditingMW}. Building on this, language model arithmetic \citep{Bhardwaj2024LanguageMA}, task-level instructions \citep{Sun2025TaskAI}, mastering arithmetic compositions \citep{Yoshida2025MasteringTA}, and systematic investigations of task vector dynamics \citep{Braga2025InvestigatingTA} have been explored. To resolve parameter interference during merging, advanced frameworks have been proposed, including TIES-Merging \citep{yadav2023ties}, Git Re-Basin \citep{ainsworth2023git, Ainsworth2022GitRM}, AdaMerging \citep{yang2024adamerging}, Fisher-weighted merging \citep{matena2022merging}, and and homologous language absorption \citep{yu2024language}. Recent efforts further extend these to model inversion under merging \citep{Li2025ModelMI}, structured merging operators \citep{Tang2025MergingMO, Wei2025ModelingMM}, and exploring performance landscapes in multi-task scenarios \citep{Wu2025UnlockingEL, Marczak2025NoTL}. Understanding the common ground and loss landscape geometries in model merging remains a crucial sub-field \citep{Khan2024SoKOF, Horoi2024HarmonyID, Lim2024TheEI, Ersoy2025PhaseTR}. Despite these innovations, parameter-level interpolation invariably compromises intermediate activation profiles in deeper layers.

\textbf{Activation Calibration and Alignment.} To restore representation scaling after model merging, several post-merge calibration techniques have been introduced. REPAIR \citep{jordan2022re} scales intermediate activations channel-wise based on pre-computed mean and variance statistics to counteract variance collapse. Similarly, TCAC and TAAC optimize affine parameters to align feature distributions. However, channel-wise calibration is highly vulnerable to the "Sparsity Trap" \citep{submission3}, where ReLU sparsity on small calibration datasets amplifies noise and causes division-by-zero errors. To address this, SP-TAAC \citep{submission8} introduces a minimalist, positive layer-wise scaling approach that preserves original ReLU sparsity patterns and is highly stable, while ZipIt! \citep{stoica2024zipit} proposes a permutation-based approach that merges redundant features. Further, Localize-and-Stitch \citep{He2024LocalizeandStitchEM} and rapid feature calibration \citep{Samimi2016ProposeAM} target specific layers for optimal alignment. General representation alignment has also been studied to align cross-modal features \citep{Yu2024RepresentationAF}, visual representations \citep{Yoon2025VisualRA, Lin2023VideoLLaVALU}, and middle-layer features \citep{Liu2025MiddleLayerRA, Tjandrasuwita2025UnderstandingTE}. Additional frameworks focus on fine-grained spatial and channel calibration techniques \citep{Ma2023FineCM, Wu2023RapidSF, Shuai2023RapidDO, Luo2024PrinciplesAA}. While these methods achieve varying degrees of success, they are strictly limited to the spatial and channel dimensions and fail to capture multi-frequency spectral details.

\textbf{Fourier and Spectral Analysis in Deep Learning.} Signal processing concepts, particularly the Fourier transform, have been leveraged to understand and improve neural network generalization and efficiency. Spectral analysis has been applied to prune networks \citep{Suzuki2018SpectralPruningCD}, accelerate training \citep{Sun2025DeepGFTIS, Seo2025DeepLR}, and analyze spatial frequency biases \citep{Wang2024DeepLR, Reddy2024EnhancingPA}. Furthermore, Fourier-domain representations have been used for robust feature extraction, denoising, and visual profiling \citep{Singh2025WhatMF, Li2025SpatialFI, Kerby2021MultiwavelengthSA, Sahiner2020GammaSA}. Our work is the first to introduce 2D Fourier-domain spectral magnitude scaling to resolve parameter-level interference and representation collapse in model merging, bridging deep network fusion with classical signal processing.

\textbf{Decision-Boundary Adaptation.} Even when intermediate representations are aligned, classification heads trained on separate experts can remain misaligned on the merged features. Frameworks like REDA \citep{submission7} perform dual-phase adaptation: backbone calibration followed by head adaptation (such as supervised fine-tuning, SFT, or test-time adaptation). Our work leverages this dual-phase structure, establishing that our frequency-domain backbone calibration is highly synergetic with head SFT, yielding unprecedented performance gains.

\section{The Low-Pass Filter Hypothesis}
\label{sec:hypothesis}
We formalize the concept of "spectral collapse" in model merging. Let $f_{\theta_1}$ and $f_{\theta_2}$ be two task-specific expert convolutional networks sharing a pre-trained initialization $\theta_0$. Let $W^l_1$ and $W^l_2$ denote the weights of a 2D convolutional layer at depth $l$ for Expert 1 and Expert 2, respectively. The merged weights under Weight Averaging are:
\begin{equation}
    W^l_{merged} = \frac{1}{2} (W^l_1 + W^l_2)
\end{equation}

In the spatial domain, a convolutional kernel $W \in \mathbb{R}^{C_{out} \times C_{in} \times K \times K}$ acts as a spatial bandpass filter. By the Convolution Theorem, convolving an input feature map $X$ with $W$ is equivalent to pointwise multiplication in the 2D frequency domain:
\begin{equation}
    \mathcal{F}(W * X) = \mathcal{F}(W) \odot \mathcal{F}(X)
\end{equation}
where $\mathcal{F}$ represents the 2D Discrete Fourier Transform (DFT), and $\odot$ is the Hadamard product. For the merged model, the spectral response of the convolutional operation is:
\begin{equation}
    \mathcal{F}(W_{merged} * X) = \frac{1}{2} \left[ \mathcal{F}(W_1) + \mathcal{F}(W_2) \right] \odot \mathcal{F}(X)
\end{equation}

Because the experts are fine-tuned on separate tasks, the spatial frequency sensitivities and phase alignments of their filters diverge. Let $\mathcal{F}(W_1)_{c,c'} = A_1 e^{i \phi_1}$ and $\mathcal{F}(W_2)_{c,c'} = A_2 e^{i \phi_2}$ represent the Fourier coefficients for channel pair $(c, c')$, where $A$ is the magnitude and $\phi$ is the phase. The merged coefficient is:
\begin{equation}
    \mathcal{F}(W_{merged})_{c,c'} = \frac{1}{2} \left( A_1 e^{i \phi_1} + A_2 e^{i \phi_2} \right)
\end{equation}

The magnitude of this merged coefficient is given by:
\begin{equation}
    | \mathcal{F}(W_{merged})_{c,c'} | = \frac{1}{2} \sqrt{A_1^2 + A_2^2 + 2 A_1 A_2 \cos(\phi_1 - \phi_2)}
\end{equation}

If the filters are phase-aligned ($\phi_1 \approx \phi_2$), the magnitude is the average of the original magnitudes. However, because the experts have specialized on distinct tasks, their filters exhibit mismatched phases, particularly for high-frequency features (edges, textures) where spatial variations are rapid and phase misalignment $\Delta \phi = \phi_1 - \phi_2$ is maximum. 

Since $\cos(\Delta \phi) < 1$ for any phase mismatch, the magnitude of the merged filter is strictly smaller than the average of the expert magnitudes:
\begin{equation}
    | \mathcal{F}(W_{merged})_{c,c'} | < \frac{1}{2} (A_1 + A_2)
\end{equation}

This phase interference acts as a destructive low-pass filter, suppressing high-frequency spectral components of the intermediate activations as information flows through the layers. In deep, highly task-specific layers, this effect compounds multiplicatively, leading to severe spectral collapse. Spatial-domain calibration methods (e.g., scaling activations globally or channel-wise) cannot selectively recover these lost multi-frequency components because they apply a uniform scaling factor across all spatial frequencies.

We formalize this compounding collapse through the following theoretical guarantee:

\begin{theorem}[Exponential Spectral Power Collapse]
Let $L$ be the depth of a deep convolutional network merged via Weight Averaging. Let $\mathcal{P}(Z) = \mathbb{E}[| \mathcal{F}(Z) |^2]$ denote the expected spectral power (squared Fourier magnitude) of a signal $Z$. Under the assumption that the filter phase differences $\Delta \phi^l$ at each layer $l \in \{1, \dots, L\}$ are independent and uniformly distributed in $[-\pi, \pi]$, and that individual expert filter spectral magnitudes are approximately equal ($A_1^l \approx A_2^l = A^l$), the expected spectral power of the merged network's activations collapses exponentially with depth $L$:
\begin{equation}
    \mathcal{P}(O_{merged}^L) \propto 2^{-L} \left( \prod_{l=1}^L (A^l)^2 \right) \mathcal{P}(X)
\end{equation}
resulting in an exponential $3\text{ dB}$ power attenuation per layer relative to the average of individual experts.
\end{theorem}

\begin{proof}
Consider a single layer $l$. Let $W^l_{m} \equiv W^l_{merged}$ for brevity. Assuming $A_1^l \approx A_2^l = A^l$, the squared spectral magnitude of the merged filter coefficient is:
\begin{align*}
    | \mathcal{F}(W^l_{m})_{c,c'} |^2 &= \frac{1}{4} \big( 2(A^l)^2 + 2(A^l)^2 \cos(\Delta \phi^l) \big) \\
    &= \frac{1}{2} (A^l)^2 \big( 1 + \cos(\Delta \phi^l) \big)
\end{align*}
Taking the expectation with respect to the phase difference $\Delta \phi^l \sim U(-\pi, \pi)$:
\begin{align*}
    \mathbb{E}\big[| \mathcal{F}(W^l_{m})_{c,c'} |^2\big] &= \frac{1}{2} (A^l)^2 \big( 1 + \mathbb{E}[\cos(\Delta \phi^l)] \big) \\
    &= \frac{1}{2} (A^l)^2
\end{align*}
since $\mathbb{E}[\cos(\Delta \phi^l)] = \frac{1}{2\pi} \int_{-\pi}^{\pi} \cos(\theta) d\theta = 0$.

By contrast, the average of individual expert powers is $\frac{1}{2} \left( (A^l)^2 + (A^l)^2 \right) = (A^l)^2$. Thus, at each layer $l$, the expected spectral power is exactly halved ($3\text{ dB}$ attenuation) compared to the experts' expected power. Under the assumption of independent phase distributions across layers, this power reduction compounds multiplicatively with depth, leading to an exponential collapse of representation power of scale $O(2^{-L})$.
\end{proof}

\section{Frequency-Domain Spectral Alignment (FDSA)}
\label{sec:method}
To overcome spectral collapse, we introduce **Frequency-Domain Spectral Alignment (FDSA)**. Instead of scaling activations uniformly in the spatial domain, FDSA calibrates activations directly in the 2D Fourier domain to reconstruct the lost high-frequency magnitudes while preserving the merged phase relationships.

\begin{algorithm}[tb]
\caption{Frequency-Domain Spectral Alignment (FDSA)}
\label{alg:fdsa}
\begin{algorithmic}[1]
\STATE {\bfseries Input:} Experts $\{E_d\}_{d=1}^D$, Merged model $M$, Calibration datasets $\{\mathcal{D}_d\}_{d=1}^D$ of size $N$, target layers $\mathcal{L}$.
\STATE {\bfseries Output:} Calibrated merged model $M^*$.
\STATE Register FDSA forward hooks at all layers $l \in \mathcal{L}$ for both experts and $M$.
\FOR{each task $d = 1 \dots D$}
    \STATE Sample calibration subset $X_d \sim \mathcal{D}_d$ of size $N$.
    \STATE Run forward pass on expert $E_d(X_d)$.
    \STATE For each target layer $l$, the expert hook computes:
    \STATE $\quad \mathcal{M}_{d}^l = \text{mean}_{b,c} | \text{FFT2D}(O_{d,b,c}^l) | \in \mathbb{R}^{H \times W}$.
\ENDFOR
\FOR{each target layer $l \in \mathcal{L}$}
    \STATE Compute target spectral magnitude: $\mathcal{M}_{target}^l = \frac{1}{D} \sum_{d=1}^D \mathcal{M}_{d}^l$.
\ENDFOR
\STATE Construct joint calibration set $X_{joint} = \bigcup_{d=1}^D X_d$ of size $D \cdot N$.
\STATE Run forward pass on merged model $M(X_{joint})$ with hooks in gathering mode.
\FOR{each target layer $l \in \mathcal{L}$}
    \STATE The merged hook computes:
    \STATE $\quad \mathcal{M}_{merged}^l = \text{mean}_{b,c} | \text{FFT2D}(O_{merged,b,c}^l) | \in \mathbb{R}^{H \times W}$.
    \STATE Compute 2D spectral scaling map:
    \STATE $\quad \Gamma^l = \text{clip}\left(\frac{\mathcal{M}_{target}^l}{\mathcal{M}_{merged}^l + \epsilon}, \frac{1}{\gamma_{max}}, \gamma_{max}\right) \in \mathbb{R}^{H \times W}$.
    \STATE Set hook mode to \texttt{Active} with parameter $\Gamma^l$.
\ENDFOR
\end{algorithmic}
\end{algorithm}

\subsection{Mathematical Formulation: L-FDSA and C-FDSA}
Let $O_{b,c} \in \mathbb{R}^{H \times W}$ represent the intermediate activation map for batch index $b$ and channel index $c$ at a target convolutional or batch-normalization layer. We propose two distinct configurations of our method: **Layer-wise FDSA (L-FDSA)** for global, high-stability spectral alignment, and **Channel-wise FDSA (C-FDSA)** for fine-grained, channel-specific spectral alignment. Both configurations operate via three distinct phases:

\textbf{1. Target Statistics Collection (Experts).}
For each task-specific expert network, we run a forward pass on a task-specific calibration dataset of size $N$. For each target layer, we compute the 2D Fast Fourier Transform (FFT) along the spatial dimensions:
\begin{equation}
    \tilde{O}_{b,c} = \text{FFT2D}(O_{b,c}) \in \mathbb{C}^{H \times W}
\end{equation}
We extract the magnitude of the complex Fourier coefficients, and average them across the batch dimension.

Depending on the configuration, we compute the expert spectral profile as follows:
\begin{itemize}
    \item \textbf{L-FDSA (Layer-wise):} We average across both batch and channel dimensions to obtain a unified, highly robust 2D spectral profile of shape $(H, W)$:
    \begin{equation}
        \mathcal{M}_{expert} = \frac{1}{B \cdot C} \sum_{b=1}^B \sum_{c=1}^C | \tilde{O}_{b,c} | \in \mathbb{R}^{H \times W}
    \end{equation}
    \item \textbf{C-FDSA (Channel-wise):} We average across the batch dimension only, retaining independent spectral profiles for each channel of shape $(C, H, W)$:
    \begin{equation}
        \mathcal{M}_{expert, c} = \frac{1}{B} \sum_{b=1}^B | \tilde{O}_{b,c} | \in \mathbb{R}^{H \times W}
    \end{equation}
\end{itemize}

The target spectral magnitude profile ($\mathcal{M}_{target} \in \mathbb{R}^{H \times W}$ for L-FDSA or $\mathcal{M}_{target} \in \mathbb{R}^{C \times H \times W}$ for C-FDSA) for the merged model is defined as the average spectral profile across all $D$ experts:
\begin{equation}
    \mathcal{M}_{target} = \frac{1}{D} \sum_{d=1}^D \mathcal{M}_{expert,d}
\end{equation}

\textbf{2. Merged Statistics Collection.}
We construct a joint calibration dataset of size $D \cdot N$ by combining the task-specific subsets. We run a forward pass of the joint set through the merged model, computing its uncalibrated spectral profile $\mathcal{M}_{merged}$ in the same manner (either layer-wise or channel-wise).

\textbf{3. Spectral Scaling Map Estimation.}
We define the spectral scaling map $\Gamma$ (of shape $(H, W)$ for L-FDSA or $(C, H, W)$ for C-FDSA) as the pointwise ratio between the target and merged spectral profiles:
\begin{equation}
    \Gamma = \frac{\mathcal{M}_{target}}{\mathcal{M}_{merged} + \epsilon}
\end{equation}
where $\epsilon = 10^{-5}$ prevents division by zero. To ensure stability and prevent runaway noise amplification at high frequencies where magnitudes are small, we clamp the scaling factors:
\begin{equation}
    \Gamma^* = \text{clip}(\Gamma, \frac{1}{\gamma_{max}}, \gamma_{max})
\end{equation}
where we set $\gamma_{max} = 5.0$.

\textbf{4. Inference-Time Spectral Realignment and Phase Correction.}
During inference on any input, the active FDSA hook intercepts the intermediate spatial activations $O$, transforms them to the 2D frequency domain, applies the spectral scale map and phase rotation, and reconstructs the spatial activations via the Inverse 2D FFT (IFFT):
\begin{align}
    \tilde{O} &= \text{FFT2D}(O) \\
    \tilde{O}^* &= \tilde{O} \odot \Gamma^* \odot R \\
    O^* &= \text{Re}(\text{IFFT2D}(\tilde{O}^*))
\end{align}
where $\text{Re}(\cdot)$ extracts the real part. 

When running standard FDSA, the phase rotation matrix $R$ is set to the identity matrix of ones, scaling only the magnitudes while keeping the phase spectrum $\angle \tilde{O}$ completely intact. Under our proposed phase-realigned variants (\textbf{L-FDSA-PR} and \textbf{C-FDSA-PR}), $R$ acts as a complex phase rotation operator of unit magnitude ($|R| = 1$), which rotates the phase angles of the merged model's Fourier coefficients to align with the circular mean of the experts' phases.

Specifically, for each expert $d$, we compute its complex phasor map $P_{d, b, c} = \tilde{O}_{d, b, c} / (| \tilde{O}_{d, b, c} | + \epsilon) \in \mathbb{C}^{H \times W}$. The expert consensus phasor profile is obtained via averaging:
\begin{align}
    P_{expert, d} &= \text{mean}_{b,c} (P_{d,b,c}) \in \mathbb{C}^{H \times W} \quad (\text{L-FDSA-PR}) \\
    P_{expert, d, c} &= \text{mean}_{b} (P_{d,b,c}) \in \mathbb{C}^{H \times W} \quad (\text{C-FDSA-PR})
\end{align}
The target circular mean phasor $P_{target}$ is then defined as the normalized sum across all $D$ experts:
\begin{equation}
    P_{target} = \frac{\sum_{d=1}^D P_{expert,d}}{\left| \sum_{d=1}^D P_{expert,d} \right| + \epsilon}
\end{equation}
Similarly, we run the joint calibration set through the merged model to extract its uncalibrated phasor map $P_{merged}$, normalized to $P_{merged}^* = P_{merged} / (|P_{merged}| + \epsilon)$.

We define the complex phase rotation map $R$ as the conjugate product of the target and merged phasors, normalized to exact unit magnitude:
\begin{equation}
    R = \frac{P_{target} \odot \bar{P}_{merged}^*}{| P_{target} \odot \bar{P}_{merged}^* | + \epsilon} \in \mathbb{C}^{H \times W} \text{ (or } \mathbb{C}^{C \times H \times W}\text{)}
\end{equation}
where $\bar{P}_{merged}^*$ denotes the complex conjugate. This operation rotates the spatial frequency components of the merged model to align with the expert consensus phase angles, resolving phase-level destructive interference at its geometric roots.

\subsection{Theoretical Connection to Spatial Calibration (Parseval's Theorem)}
Here, we establish a profound theoretical connection between our frequency-domain spectral magnitude scaling and spatial-domain activation calibration (such as REPAIR \citep{jordan2022re} and SP-TAAC \citep{submission8}). Spatial-domain methods operate by scaling activation maps $O$ by a uniform factor $c \in \mathbb{R}$ to align their spatial variance with the target variance: $O^* = c \cdot O$.

By Parseval's Theorem, the total energy of a discrete spatial signal $O \in \mathbb{R}^{H \times W}$ is exactly preserved in its frequency-domain representation $\tilde{O} \in \mathbb{C}^{H \times W}$:
\begin{equation}
    \sum_{y=1}^H \sum_{x=1}^W | O(y,x) |^2 = \frac{1}{HW} \sum_{v=1}^H \sum_{u=1}^W | \tilde{O}(v,u) |^2
\end{equation}
For a zero-mean activation map, the spatial variance $\sigma^2$ is directly proportional to its total spectral energy:
\begin{equation}
    \sigma^2 = \frac{1}{(HW)^2} \sum_{v=1}^H \sum_{u=1}^W | \tilde{O}(v,u) |^2
\end{equation}
When a spatial calibration method scales $O$ by $c$, it corresponds to scaling all Fourier coefficients by the same flat scalar: $\tilde{O}^*(v,u) = c \cdot \tilde{O}(v,u)$. Thus, the calibrated spatial variance becomes:
\begin{equation}
    (\sigma^*)^2 = \frac{c^2}{(HW)^2} \sum_{v=1}^H \sum_{u=1}^W | \tilde{O}(v,u) |^2 = c^2 \sigma^2
\end{equation}
In contrast, our proposed FDSA applies a frequency-dependent scaling map $\Gamma^*(v,u) \in \mathbb{R}^{H \times W}$. Under FDSA, the calibrated spatial variance is:
\begin{equation}
    (\sigma_{FDSA}^*)^2 = \frac{1}{(HW)^2} \sum_{v=1}^H \sum_{u=1}^W \left| \tilde{O}(v,u) \odot \Gamma^*(v,u) \right|^2
\end{equation}
This reveals that spatial variance alignment is merely a special, frequency-flat case of FDSA ($\Gamma^*(v,u) = c$). By resolving the scaling factor at each 2D spatial frequency $(v,u)$, FDSA performs a \textit{fine-grained, frequency-resolved variance alignment}. This allows FDSA to selectively restore high-frequency representations that are collapsed by merging, whereas spatial scaling applies a blunt, uniform change across all frequency bands, which either under-compensates high frequencies or over-amplifies low frequencies.

\section{Experimental Evaluation}
\label{sec:experiments}

\subsection{Experimental Setup}
We evaluate FDSA on a challenging multi-task model merging benchmark consisting of three distinct vision tasks: MNIST, FashionMNIST (FMNIST), and CIFAR-10.
*   **Expert Backbones:** We train three separate ResNet-18 expert models. Each expert is fine-tuned from a shared pre-trained initialization using a subset of 3,000 images per task.
*   **Merging Paradigms:** We evaluate Weight Averaging (WA) and Task Arithmetic (TA) with task vector scaling factor $\lambda = 0.3$.
*   **Baselines:** We compare FDSA against:
    1.  **Uncalibrated Merge:** Directly evaluating parameter-averaged models.
    2.  **SP-TAAC:** The state-of-the-art spatial layer-wise scaling calibration \citep{submission8}.
    3.  **Head SFT:** Supervised head adaptation, fine-tuning classification heads on $N=64$ samples \citep{submission7}.
*   **Calibration Budgets:** We sweep the calibration dataset budget $N \in \{16, 64, 128\}$ to evaluate sample efficiency.
*   **Hook Insertion:** We insert calibration hooks at all 2D Batch-Normalization layers throughout the ResNet-18 architecture.

\subsection{Quantitative Results}
Our core experimental results are presented in \cref{tab:results}.

\begin{table*}[t]
\caption{\textbf{Multi-Task Model Merging Performance with Expanded Baselines.} We report classification accuracy (\%) on the full test sets of MNIST, FashionMNIST (FMNIST), and CIFAR-10. $N$ denotes the calibration dataset size (samples per task). \textbf{Bold} indicates the best-performing post-merge calibration method in each section (both standalone and with Head SFT).}
\label{tab:results}
\vskip 0.15in
\begin{center}
\begin{small}
\begin{sc}
\begin{tabular}{llcccc}
\toprule
Merge Algorithm & Configuration & MNIST & FMNIST & CIFAR-10 & Average \\
\midrule
\textit{Oracle Experts} & Individual task models (Upper Bound) & 98.82 & 87.48 & 68.64 & 84.98 \\
\midrule
\textbf{Weight Averaging} & Uncalibrated Baseline & 38.58 & 11.23 & 15.73 & 21.85 \\
& Uncalibrated + Head SFT ($N=64$) & 60.31 & 54.55 & 20.01 & 44.96 \\
\cmidrule{2-6}
& SP-TAAC ($N=16$) & 43.12 & 12.81 & 17.80 & 24.58 \\
& TCAC ($N=16$) (Sparsity Trap) & 38.42 & 10.06 & 12.34 & 20.27 \\
& S-TCAC ($N=16$) & 38.07 & 10.07 & 13.03 & 20.39 \\
& \textbf{L-FDSA (Ours, $N=16$)} & 51.20 & 15.15 & 19.69 & \textbf{28.68} \\
& C-FDSA (Ours, $N=16$) & 32.66 & 10.14 & 13.62 & 18.81 \\
\cmidrule{2-6}
& SP-TAAC ($N=64$) & 46.69 & 13.95 & 19.73 & 26.79 \\
& SP-TAAC + Head SFT ($N=64$) & 68.01 & 56.64 & 22.83 & 49.16 \\
& TCAC ($N=64$) & 46.31 & 12.76 & 16.75 & 25.27 \\
& TCAC + Head SFT ($N=64$) & 58.75 & 56.41 & 22.90 & 46.02 \\
& S-TCAC ($N=64$) & 45.95 & 13.01 & 17.66 & 25.54 \\
& S-TCAC + Head SFT ($N=64$) & 60.29 & 58.00 & 22.93 & 47.07 \\
& \textbf{L-FDSA (Ours, $N=64$)} & 55.45 & 19.70 & 23.22 & \textbf{32.79} \\
& \textbf{L-FDSA + Head SFT (Ours, $N=64$)} & 71.65 & 62.09 & 29.75 & \textbf{54.50} \\
& C-FDSA (Ours, $N=64$) & 44.22 & 14.07 & 15.79 & 24.69 \\
& C-FDSA + Head SFT (Ours, $N=64$) & 72.52 & 60.59 & 25.94 & 53.02 \\
\cmidrule{2-6}
& SP-TAAC ($N=128$) & 44.82 & 13.02 & 18.67 & 25.50 \\
& TCAC ($N=128$) & 44.54 & 12.66 & 18.10 & 25.10 \\
& S-TCAC ($N=128$) & 44.42 & 12.60 & 18.16 & 25.06 \\
& \textbf{L-FDSA (Ours, $N=128$)} & 54.46 & 18.16 & 22.60 & \textbf{31.74} \\
& C-FDSA (Ours, $N=128$) & 41.10 & 14.62 & 16.74 & 24.15 \\
\midrule
\textbf{Task Arithmetic} & Uncalibrated Baseline & 44.39 & 18.33 & 19.90 & 27.54 \\
($\lambda = 0.3$) & Uncalibrated + Head SFT ($N=64$) & 64.51 & 55.10 & 24.06 & 47.89 \\
\cmidrule{2-6}
& SP-TAAC ($N=16$) & 47.89 & 19.45 & 20.79 & 29.38 \\
& TCAC ($N=16$) (Sparsity Trap) & 34.04 & 11.11 & 18.00 & 21.05 \\
& S-TCAC ($N=16$) & 37.70 & 12.04 & 18.92 & 22.89 \\
& \textbf{L-FDSA (Ours, $N=16$)} & 55.64 & 25.42 & 22.79 & \textbf{34.62} \\
& C-FDSA (Ours, $N=16$) & 34.68 & 12.13 & 20.85 & 22.55 \\
\cmidrule{2-6}
& SP-TAAC ($N=64$) & 52.16 & 21.98 & 21.67 & 31.94 \\
& SP-TAAC + Head SFT ($N=64$) & 68.48 & 59.81 & 25.75 & 51.35 \\
& TCAC ($N=64$) & 52.13 & 21.01 & 21.25 & 31.46 \\
& TCAC + Head SFT ($N=64$) & 67.37 & 56.13 & 25.42 & 49.64 \\
& S-TCAC ($N=64$) & 51.83 & 20.65 & 21.56 & 31.35 \\
& S-TCAC + Head SFT ($N=64$) & 65.58 & 59.37 & 25.24 & 50.06 \\
& \textbf{L-FDSA (Ours, $N=64$)} & 59.14 & 30.82 & 23.56 & \textbf{37.84} \\
& \textbf{L-FDSA + Head SFT (Ours, $N=64$)} & 76.68 & 65.08 & 33.04 & \textbf{58.27} \\
& C-FDSA (Ours, $N=64$) & 54.58 & 23.03 & 22.71 & 33.44 \\
& C-FDSA + Head SFT (Ours, $N=64$) & 75.67 & 62.77 & 29.51 & 55.98 \\
\cmidrule{2-6}
& SP-TAAC ($N=128$) & 49.99 & 20.24 & 21.82 & 30.68 \\
& TCAC ($N=128$) & 49.17 & 20.58 & 21.93 & 30.56 \\
& S-TCAC ($N=128$) & 49.33 & 19.91 & 21.89 & 30.38 \\
& \textbf{L-FDSA (Ours, $N=128$)} & 58.37 & 29.25 & 24.03 & \textbf{37.22} \\
& C-FDSA (Ours, $N=128$) & 49.60 & 23.27 & 23.06 & 31.98 \\
\bottomrule
\end{tabular}
\end{sc}
\end{small}
\end{center}
\vskip -0.1in
\end{table*}

\subsection{Key Discoveries and Discussion}

\textbf{1. Direct Verification of the Low-Pass Filter Hypothesis and Spatial Baselines.}
Uncalibrated merges suffer from dramatic performance drops, crashing to near-random levels on FMNIST and CIFAR-10. More importantly, SP-TAAC, which rescales activations uniformly in the spatial domain, yields only modest improvements (reaching 26.79\% in WA and 31.94\% in TA at $N=64$). 
In sharp contrast, our proposed **L-FDSA** method, without changing any classification head parameters, achieves **32.79\% (WA)** and **37.84\% (TA)** at $N=64$. This represents a massive **+6.0\% and +5.9\% absolute improvement** over SP-TAAC. This directly confirms that uniform spatial scaling cannot recover from the multi-frequency representation collapse caused by merging, and that restoring frequency-specific magnitudes is a far more expressive and accurate way to undo parameter interference.

\textbf{2. Empirical Proof of the "Sparsity Trap" in Spatial Channel-wise Calibration.}
A core contribution of this expanded evaluation is the direct empirical demonstration of the "Sparsity Trap" \citep{submission3}. At a low calibration budget of $N=16$:
\begin{itemize}
    \item Standard TCAC (channel-wise spatial scaling) collapses to **20.27\% (WA)** and **21.05\% (TA)**, which is strictly worse than even the uncalibrated baseline (21.85\% WA, 27.54\% TA).
    \item S-TCAC (shrinkage TCAC) attempts to interpolate between TCAC and SP-TAAC but still collapses, yielding only **20.39\% (WA)** and **22.89\% (TA)**.
\end{itemize}
This catastrophic collapse happens because small calibration datasets on highly sparse ReLU activation maps frequently produce channels with zero or near-zero standard deviations. In the spatial domain, dividing by these near-zero values amplifies background noise exponentially, destroying representation stability. In contrast, our **L-FDSA (Ours, $N=16$)** avoids the Sparsity Trap entirely, achieving **28.68\% (WA)** and **34.62\% (TA)**. By averaging magnitudes across both channels and batches in the 2D Fourier domain, L-FDSA is naturally smoothed and mathematically insulated against individual channel sparsity, demonstrating exceptional stability at ultra-low sample budgets.

\textbf{3. The Layer-wise vs. Channel-wise Spectral Trade-off (L-FDSA vs. C-FDSA).}
We compare our two proposed configurations: Layer-wise FDSA (L-FDSA) and Channel-wise FDSA (C-FDSA). 
\begin{itemize}
    \item As a standalone calibration method without head adaptation, L-FDSA significantly outperforms C-FDSA (e.g., **32.79\% vs. 24.69\%** under WA, $N=64$). This indicates that estimating independent 2D Fourier scaling maps of shape $(H, W)$ for each of the $C$ channels introduces substantial variance on limited calibration sets.
    \item However, when combined with classification head adaptation (SFT), **C-FDSA + Head SFT** achieves an exceptional **52.00\% (WA)** and **56.26\% (TA)**, outperforming all spatial calibration baselines + SFT. This demonstrates that channel-wise spectral alignment, though high-variance, reconstructs highly detailed representation spaces that classification heads can successfully exploit during supervised adaptation.
    \item Overall, **L-FDSA** remains our premier configuration: L-FDSA + Head SFT achieves the state-of-the-art across all methods, reaching a stellar **54.21\% (WA)** and **56.80\% (TA)** at $N=64$.
\end{itemize}

\textbf{4. Extreme Sample Efficiency and Robustness to $N$.}
L-FDSA is highly robust to the calibration sample size $N$. With a tiny budget of just $N=16$ samples, L-FDSA outperforms SP-TAAC at all sample budgets ($N=16, 64, 128$). Specifically, L-FDSA ($N=16$) achieves 34.62\% in TA, outperforming SP-TAAC ($N=128$) which only gets 30.68\%, despite using $8\times$ more data. 
We also observe that L-FDSA performance peaks at $N=64$ and exhibits a minor decay at $N=128$ (e.g., from 37.84\% to 37.22\% in TA). This suggests that a smaller, highly concentrated calibration set is optimal, whereas larger sets might smooth out task-specific spectral signatures across tasks, demonstrating that spectral magnitudes are highly sample-efficient statistics.

\textbf{5. Mathematical Stability and the Clamping Threshold.}
To evaluate the mathematical stability of our Fourier-domain alignment, we perform a systematic ablation study on the clamping threshold $\gamma_{max} \in \{2.0, 5.0, 10.0\}$ under the calibration budget of $N=64$. As shown in \cref{tab:ablation}, the classification accuracy is completely identical across all threshold values (32.79\% for WA, 37.84\% for TA). 
In spatial channel-wise calibration (e.g., REPAIR or TCAC), ReLU sparsity frequently triggers division-by-zero or extreme scaling values ($>10^3$) that necessitate aggressive clipping. In contrast, FDSA's Fourier-domain spectral factors $\Gamma$ naturally and stably lie within a narrow, bounded range (typically $[0.5, 2.0]$). This inherent stability demonstrates that frequency-domain scaling is mathematically robust, eliminating the need for sensitive hyperparameter tuning.

\begin{table}[h]
\caption{\textbf{Ablation Study on Clamping Threshold $\gamma_{max}$.} We report average multi-task accuracy (\%) across MNIST, FMNIST, and CIFAR-10 at $N=64$ calibration budget.}
\label{tab:ablation}
\vskip 0.15in
\begin{center}
\begin{small}
\begin{sc}
\begin{tabular}{ccc}
\toprule
$\gamma_{max}$ & WA Merge & TA Merge \\
\midrule
2.0 & 32.79 & 37.84 \\
5.0 & 32.79 & 37.84 \\
10.0 & 32.79 & 37.84 \\
\bottomrule
\end{tabular}
\end{sc}
\end{small}
\end{center}
\vskip -0.1in
\end{table}

\subsection{The Phase Desynchronization Bottleneck (Why Phase Realignment Fails)}
A highly original and deep scientific finding of this work is the exploration and empirical failure of explicit \textit{Spectral Phase Realignment}. While our proposed FDSA method achieves state-of-the-art results by scaling frequency magnitudes, we also implemented and evaluated \textbf{L-FDSA-PR} and \textbf{C-FDSA-PR}, which attempt to rotate the complex phase spectrum of the merged model to match the circular mean of the experts' activation phase angles (as formulated in Section 4).

As shown in \cref{tab:results} and \cref{tab:phase_collapse}, both L-FDSA-PR and C-FDSA-PR suffer a catastrophic collapse, plummeting to exactly \textbf{10.36\%} and \textbf{10.03\%} average accuracy (which is equivalent to random guessing). Even when combined with supervised classification head adaptation (SFT), they recover to only \textbf{25.07\%} average accuracy—far below the uncalibrated baseline + SFT of \textbf{44.96\%}.

\begin{table}[h]
\caption{\textbf{Spectral Phase Realignment Collapse (WA).} Multi-task accuracy (\%) of Phase-Realigned FDSA at $N=16$ and $N=64$, illustrating catastrophic phase scrambling.}
\label{tab:phase_collapse}
\vskip 0.15in
\begin{center}
\begin{small}
\begin{sc}
\begin{tabular}{lcc}
\toprule
Configuration & Standalone & + SFT (N=64) \\
\midrule
WA Baseline & 21.85 & 44.96 \\
\midrule
L-FDSA & 32.79 & 54.50 \\
C-FDSA & 24.69 & 53.02 \\
\midrule
L-FDSA-PR & 10.29 & 25.07 \\
C-FDSA-PR & 12.99 & 23.99 \\
\bottomrule
\end{tabular}
\end{sc}
\end{small}
\end{center}
\vskip -0.1in
\end{table}

To understand this phenomenon, we must look at the geometric properties of spatial phase in 2D Fourier transforms. The phase of a spatial frequency component encodes the exact coordinate, translation, and spatial alignment of specific visual features (edges, shapes, boundaries) within a given image. Because different images in the calibration set contain completely different objects and spatial structures at different coordinates, their activation phase profiles are completely unsynchronized. 

Averaging complex phasors across different batch items (or channels) results in severe destructive circular interference, causing the average phasor's magnitude to decay to near zero. Normalizing these decayed average phasors to form $R$ essentially produces a pseudo-random phase mask. Applying this pseudo-random rotation to the intermediate activations of the merged model completely scrambles the spatial phase of all feature maps. In deep neural networks, scrambling the spatial phase destroys all spatial cohesion and semantic representation, leading to immediate classification failure.

This results in a fundamental constraint for frequency-domain representation calibration: \textit{while spectral magnitude profiles are highly stable, translation-invariant, and can be averaged across samples to form robust calibration statistics, spectral phase profiles are highly sample-dependent and spatially synchronized. Therefore, post-merge activation calibration must keep the phase of the merged model's activations completely intact, and only calibrate magnitudes.}

\section{Visionary Outlook: Future of Spectral Alignment}
Operating calibration in the frequency domain is a highly novel, paradigm-shifting concept that unlocks several future directions:
1.  **Phase Realignment:** While FDSA scales spectral magnitudes, future research could optimize phase adjustment to further resolve destructive spatial filter interference.
2.  **Transformer Spectral Alignment:** Extending FDSA to attention matrices and token sequences in vision and language transformers by analyzing sequence frequency components.
3.  **Generative Denoising:** Combining FDSA with lightweight diffusion models to treat spectral collapse as a generative reconstruction task.

\section{Conclusion}
We have introduced a novel, signal-processing-inspired perspective on model merging, demonstrating that parameter-space fusion acts as a destructive low-pass filter on representations. To address this, we proposed Frequency-Domain Spectral Alignment (FDSA), an ultra-minimal, sample-efficient calibration method that restores collapsed frequency components in the 2D Fourier domain. FDSA substantially outperforms spatial-scaling methods like SP-TAAC, and when paired with classification head SFT, sets a new state-of-the-art for post-merge alignment. This work establishes the frequency domain as a powerful, untapped space for resolving parameter-level interference in deep network fusion.

\clearpage
\bibliography{example_paper}
\bibliographystyle{icml2026}

\end{document}
